\documentclass[a4paper,10pt]{report}
\usepackage{graphicx}
\usepackage[utf8]{inputenc}
\usepackage[margin=1.5cm]{geometry} % Reduced margins for less whitespace
\usepackage{multicol}
\usepackage{amsmath}
\usepackage{array}
\usepackage{pgfplots}
\usepackage{palatino}
\usepackage{fancyhdr}
\usepackage{setspace}
\usepackage{listings}
\usepackage{tikz}

\pgfplotsset{compat=1.18}

\lstset{frame=, basicstyle={\footnotesize\ttfamily}}

\begin{document}

\raggedbottom
\setstretch{1.05}

\pagestyle{fancy}
\fancyhf{}
\fancyhead[L]{\small Elementary Charge and Millikan Experiment\par}
\fancyfoot[C]{\thepage}
\setlength{\headheight}{14pt}

% Compact Title Section
\begin{center}
{\Large \bfseries Elementary Charge and Millikan Experiment\par}
\vspace{3pt}
{\normalsize Muhammed Emin Duran | 0412220007\par}
\vspace{4pt}
\rule{\textwidth}{1pt}
\end{center}

\vspace{2pt}

\setlength{\columnsep}{20pt}
\setlength{\columnseprule}{0.5pt}

\begin{multicols}{2}

\section*{Introduction}
The purpose of this experiment is to determine the elementary charge $e$ using charged
oil droplets in the Millikan experiment. From their rise and fall times at different
voltages, we calculate the droplets' charges and radii, then determine $e$ from their
velocities in and against the direction of gravity. Finally, we estimate the electron's
$e/m$ ratio.
 
\section*{Theory and Methods}
The rising and falling motion of charged oil droplets in the electric field of the
capacitor is observed.\\

The force $F$ applied to a sphere of radius $r$ and velocity $v$ in a viscous liquid of
viscosity $\eta$ is given by Stokes' law:
\begin{equation}
    F_d = 6\pi r \eta v 
\end{equation}

A droplet of mass $m$, volume $V$, and density $\rho_1$ is located in Earth's
gravitational field.
\begin{equation}
    F_g = mg = \rho_1 V g
\end{equation}  

Buoyant force $F_b$ is given by:
\begin{equation}
    F_b = \rho_2 V g, \quad \rho_2 \text{ :is the density of air}
\end{equation}

Electric force $F_e$ is given by:
\begin{equation}
    F_e = QE, \quad E = \frac{U}{d}
\end{equation}

While $Q$ is the charge on the droplets, ($Q = ne$) $U$ is the applied voltage,
and $d$ is the distance between the plates of the capacitor.\\

\begin{center}
    \centering
    \includegraphics[width=\linewidth]{figures/ForceDiagram.png}
    \\
    \textbf{Figure:} Forces acting on a droplet in the Millikan experiment.
\end{center}

\textbf{1. Falling Droplet ($v_1$)} \\
When the droplet is falling, the forces acting on it are:
\begin{equation}
    F_b + F_d = F_g + F_e
\end{equation}
\begin{gather}
    \frac{4}{3}\pi r^3 \rho_2 g + 6\pi \eta r v_1 = \frac{4}{3}\pi r^3 \rho_1 g + Q\frac{U}{d} \\
    6\pi \eta r v_1 = Q\frac{U}{d} + \frac{4}{3}\pi r^3 g (\rho_1 - \rho_2) \\
    v_1 = \frac{1}{6\pi \eta r} \left( Q\frac{U}{d} + \frac{4}{3}\pi r^3 g (\rho_1 - \rho_2) \right)
\end{gather}

\textbf{2. Rising Droplet ($v_2$)} \\
When the droplet is rising, the forces acting on it are:
\begin{equation}
    F_b + F_e = F_g + F_d
\end{equation}
\begin{gather}
    \frac{4}{3}\pi r^3 \rho_2 g + Q\frac{U}{d} = \frac{4}{3}\pi r^3 \rho_1 g + 6\pi \eta r v_2 \\
    6\pi \eta r v_2 = Q\frac{U}{d} - \frac{4}{3}\pi r^3 g (\rho_1 - \rho_2) \\
    v_2 = \frac{1}{6\pi \eta r} \left( Q\frac{U}{d} - \frac{4}{3}\pi r^3 g (\rho_1 - \rho_2) \right)
\end{gather}

By subtracting the equation for $v_2$ from $v_1$, the electric force terms cancel out, allowing us to 
isolate the radius $r$:
\begin{gather*}
    v_1 - v_2 = \frac{4 r^2 g (\rho_1 - \rho_2)}{9\eta} \\
    r^2 = \frac{9\eta (v_1 - v_2)}{4g (\rho_1 - \rho_2)} \\
    r = \frac{3}{2} \sqrt{\frac{\eta}{g(\rho_1 - \rho_2)}} \sqrt{|v_1 - v_2|}
\end{gather*}
By grouping the experimental constants into $C_2$, the equation for the radius simplifies
to:
\begin{gather}
    C_2 = 6.37 \cdot 10^{-5} \, (\text{m} \cdot \text{s})^{1/2} \\
    r = C_2 \sqrt{|v_1 - v_2|}
\end{gather}

Similarly, by adding the equations for $v_1$ and $v_2$, the gravitational and buoyant 
force terms cancel out, allowing us to find the charge $Q$:
\begin{gather*}
    v_1 + v_2 = \frac{1}{6\pi \eta r} \left( 2Q\frac{U}{d} \right) = \frac{QU}{3\pi \eta r d} \\
    Q = \frac{3\pi \eta d (v_1 + v_2) r}{U}
\end{gather*}
Substituting the full expression for $r$ derived above into this equation gives:
\begin{gather*}
    Q = \left( \frac{9}{2} \pi d \sqrt{\frac{\eta^3}{g(\rho_1 - \rho_2)}} \right) \frac{v_1 + v_2}{U} \sqrt{|v_1 - v_2|}
\end{gather*}
By grouping these constants into $C_1$, the final equation for the charge of the droplet 
is obtained:
\begin{gather}
    C_1 = 2.73 \cdot 10^{-11} \, \text{kg} \cdot \text{m} (\text{m} \cdot \text{s})^{-1/2} \\
    Q = C_1 \frac{v_1 + v_2}{U} \sqrt{|v_1 - v_2|}
\end{gather}


\subsection*{Experimental Procedure}
The eyepiece micrometer was calibrated using a stage micrometer. Oil droplets were
introduced with the sprayer, and the microscope was focused on a charged droplet. Using
the commutator switch, the droplet was moved between two calibrated positions while its
ascent and descent times were recorded. The measurements were repeated at 400 V and
500 V. The distances and times were then used to calculate the velocities, charges, and
radii of the droplets.

\section*{Results}
\begin{center}
\textbf{Table 1: Measured ascent and descent data} \\
\vspace{5pt}
\renewcommand{\arraystretch}{1.2}
\resizebox{\columnwidth}{!}{%
\begin{tabular}{c c c c c c c}
\hline
$U$ (V) & $t_1$ (s) & $s_1$ (div) & $s_1$ (mm) & $t_2$ (s) & $s_2$ (div) & $s_2$ (mm) \\
\hline
400 & 3.23 & 30 & 0.89 & 3.25 & 30 & 0.89 \\
400 & 3.43 & 30 & 0.89 & 3.06 & 30 & 0.89 \\
400 & 1.65 & 30 & 0.89 & 2.33 & 30 & 0.89 \\
500 & 3.23 & 30 & 0.89 & 3.53 & 30 & 0.89 \\
500 & 2.12 & 30 & 0.89 & 2.82 & 30 & 0.89 \\
\hline
\end{tabular}%
}
\end{center}

\vspace{10pt}
\begin{center}
\textbf{Table 2: Calculated velocities, charge, and radius} \\
\vspace{5pt}
\renewcommand{\arraystretch}{1.2}
\resizebox{\columnwidth}{!}{%
\begin{tabular}{c c c c c}
\hline
$U$ (V) & $v_1$ falling (m/s) & $v_2$ rising (m/s) & $Q$ (C) & $r$ (m) \\
\hline
400 & $2.7\times10^{-4}$ & $2.8\times10^{-4}$ & $4.9\times10^{-20}$ & $8.3\times10^{-8}$ \\
400 & $2.9\times10^{-4}$ & $2.6\times10^{-4}$ & $2.1\times10^{-19}$ & $3.6\times10^{-7}$ \\
400 & $3.8\times10^{-4}$ & $5.4\times10^{-4}$ & $7.9\times10^{-19}$ & $8.0\times10^{-7}$ \\
500 & $2.5\times10^{-4}$ & $2.8\times10^{-4}$ & $1.4\times10^{-19}$ & $3.1\times10^{-7}$ \\
500 & $3.2\times10^{-4}$ & $4.2\times10^{-4}$ & $4.1\times10^{-19}$ & $6.5\times10^{-7}$ \\
\multicolumn{3}{r}{\textbf{Average } $Q$:} & $3.2\times10^{-19}$ & -- \\
\hline
\end{tabular}%
}
\end{center}

\noindent\textbf{Sample calculations:} For the first 400 V measurement,
\begin{gather*}
v_1=\frac{0.00089}{3.25}=2.73846\times10^{-4}\ \mathrm{m/s},\\
v_2=\frac{0.00089}{3.23}=2.75542\times10^{-4}\ \mathrm{m/s}.
\end{gather*}
Using the manual constants $C_1=2.73\times10^{-11}$ and
$C_2=6.37\times10^{-5}$,
\begin{align*}
Q &= (2.73\times10^{-11})
\frac{(2.73846\times10^{-4})+(2.75542\times10^{-4})}{400}\\
&\quad\times\sqrt{|2.73846\times10^{-4}-2.75542\times10^{-4}|}\\
&=4.8826\times10^{-20}\ \mathrm{C},\\
r &= (6.37\times10^{-5})
\sqrt{|2.73846\times10^{-4}-2.75542\times10^{-4}|}\\
&=8.2948\times10^{-8}\ \mathrm{m}.
\end{align*}

For the first 500 V measurement, $v_1=0.00089/3.53=2.5\times10^{-4}$ m/s
and $v_2=0.00089/3.23=2.8\times10^{-4}$ m/s. Therefore,
$Q=1.4\times10^{-19}$ C and $r=3.1\times10^{-7}$ m.

\begin{center}
\includegraphics[width=\columnwidth]{figures/plots/charge_vs_radius.png}
\textbf{Figure:} Measured charge $Q$ as a function of droplet radius $r$.
\end{center}

\begin{center}
\textbf{Table 3: Results} \\
\vspace{5pt}
\begin{tabular}{l c}
\hline
\textbf{Investigated levels} & $\Delta Q$ (C) \\
\hline
$Q=0$ to first level & $4.9\times10^{-20}$ \\
First to second level & $9.1\times10^{-20}$ \\
Second to third level & $7.0\times10^{-20}$ \\
Third to fourth level & $2.0\times10^{-19}$ \\
Fourth to fifth level & $3.8\times10^{-19}$ \\
\hline
\textbf{Average} & $1.580\times10^{-19}$ \\
\hline
\end{tabular}
\end{center}

The elementary charge was estimated from the average charge-level difference:
\[
e_{\mathrm{exp}}=1.580\times10^{-19}\ \mathrm{C}.
\]
This is close to the literature value
$e_{\mathrm{lit}}=1.602\times10^{-19}\ \mathrm{C}$.

\section*{Discussion}

The measured elementary charge is close to the accepted value. The percentage error is
\[
\frac{|1.580-1.602|}{1.602}\times100=1.37\%.
\]
The small difference may come from errors in measuring time and distance, focusing the
microscope, and reading the micrometer. The graph shows different measured charge
values. Their average spacing gives an estimate of the elementary charge.

Only five droplets were measured, so more measurements would give a stronger result.
The electron mass was not measured; therefore, $e/m_e$ was calculated using the accepted
mass:
\[
\frac{e}{m_e}=\frac{1.580\times10^{-19}}{9.109\times10^{-31}}
=1.734\times10^{11}\ \mathrm{C/kg}.
\]

Overall, the measurements at 400 V and 500 V gave
$e=1.580\times10^{-19}\ \mathrm{C}$ and
$e/m_e=1.734\times10^{11}\ \mathrm{C/kg}$. The $e/m_e$ value is derived because
the electron mass was taken from the accepted value.



\end{multicols}
\end{document}